% We outline the task specification, define key terms, describe the extraction and annotation process, and present the evaluation criteria (Figure \ref{fig:flowchart}). 
We next describe our dataset construction and annotation pipeline, define the forecasting and survival prediction tasks, detail the modeling approaches evaluated, and present the sensitivity analyses used to assess robustness (Figure~\ref{fig:flowchart}).

\begin{figure}
    \centering
    \includegraphics[width=\linewidth]{AAAI_2026/Figures/ttsforecast_flowchart.png}
    \caption{T2S2 forecasting analysis pipeline}
    \label{fig:flowchart}
    \vspace{-5mm}
\end{figure}


\subsection{Dataset Extraction}

We utilized the PubMed Open Access (PMOA) Subset (as of April 25, 2024) for our analysis. Building on the method described in prior work \cite{wang2025large}, the body of each case was extracted by selecting the text between the markers \texttt{==== Body} and \texttt{==== Ref}, which serve as section delimiters in the PMOA corpus. To maintain relevance, we include only documents containing case-insensitive matches to \texttt{(case report|case presenta)} and \texttt{year-? ?old}, an approach shown to have higher specificity than PubMed metadata \cite{noroozizadeh2025reconstructing}.


As demonstrated in prior work of \cite{noroozizadeh2025reconstructing}, additional filtering was applied to identify potential sepsis-related case reports using the regular expressions \texttt{sepsi|\allowbreak septic} and \texttt{critical|intensive} care. We then used LLM-based queries to extract key attributes, including sepsis diagnosis, patient count, age, and gender. For this step, both \textit{GPT-4o} and \textit{Llama-3.1-8B-Instruct} were employed. Reports describing more than one case were excluded. A case report was retained if either model classified it as involving a Sepsis-3 diagnosis.


This multi-step process yielded a total of 2319 Sepsis-3 case reports. From this set, we randomly selected two subsets for evaluation: a group of 10 reports (sepsis-10), which underwent expert clinical annotation to serve as a gold standard, and a larger group of 100 reports (sepsis-100) as a bronze standard for broader testing based on an alternative annotator (O1-preview; Figure \ref{fig:flowchart}). %We split this data into training (n = ), validation (256), and test (244) sets.

\subsection{Textual Time Series Annotation}
\label{sec:tts-annotation}

Following \citeauthor{noroozizadeh2025reconstructing} \citeyearpar{noroozizadeh2025reconstructing}, a textual time series case report corpus of Sepsis-3 patients was constructed from the PMOA sepsis reports. A textual time series refers to a sequence of clinical findings, each paired with a timestamp (either absolute or relative to the case presentation time), corresponding to an individual patient. Here, a clinical finding is defined as a free-text expression describing an entity that pertains to or may impact a person's health.

A textual time series differs from conventional approaches to clinical concept annotation (e.g., \cite{sun2013evaluating,uzuner20112010}) in its extension of the text span to better capture the specificity and contextual meaning of each clinical finding. In our work, the interpretation of a clinical finding differs from the i2b2 concept guidelines in the following respects : (1) Clinical findings are not restricted to single prepositional phrases following a markable entity. For instance, instead of splitting “pain in chest that radiates substernally” into less informative parts, we retain the full phrase as a single, complete finding to preserve its meaning and context. (2) To improve clarity, coordinated phrases are broken into individual findings. For example, “metastases in the liver and pancreas” becomes “metastasis in the liver” and “metastasis in the pancreas.”

We used \textit{DeepSeek-R1-UD-IQ1} and \textit{Llama-3.3-70B-Instruct} models to generate clinical textual time series from the 2319 sepsis case reports. We refer to this as \textbf{T}extual \textbf{T}ime-\textbf{S}eries for \textbf{S}epsis (T2S2). The exact LLM extraction prompt can be found in Appendix~\ref{TTS_prompt}. Our dataset was split into training (n = 1,795; 80\% train, 20\% validation), testing T2S2-test (n = 244), and two external validation sets: sepsis-10 (n = 10) and sepsis-100 (n = 90).


%%%%%%%%%%%%%%%%%%%%%%%%%%%%%%%%%%%%%%%%%%%%%%%%%%%%%%%%%%%%%%%
%%%%%%%%%%%%%%%%%%%%%%%%%%%%%%%%%%%%%%%%%%%%%%%%%%%%%%%%%%%%%%
\subsection{Forecasting Tasks}
Using the timelines extracted from the annotations, we defined two primary forecasting tasks and a survival analysis task as follows:

\paragraph{Event Occurrence Prediction.} 
Given a prefix of the clinical timeline (all events up to a certain time $t$), the model is tasked with predicting whether each of the immediate next $k$ events occurs within a specified time horizon. 
This setup is repeated across multiple time cut-offs to simulate an “online” forecasting scenario, where the model must output a binary label for each of the next $k$ events:
does this event occur within $h$ hours after $t$? 
Time horizons used include 1 hour, 24 hours (1 day), and 168 hours (1 week). 
The task is framed as a series of binary classification problems, with one binary decision per event. 
Evaluation is based on precision, recall, and F1 score, averaged across event positions for each time horizon.

\paragraph{Temporal Ordering Prediction.} 
This task assesses the model's ability to reconstruct the correct sequence of future events. 
At each time cut-off $t$, we extract the next $k$ events from the timeline and remove their timestamps. 
The model must correctly output a permutation over these events that matches their true chronological order. 
This is framed as a ranking task, evaluated by computing the pairwise concordance between the predicted and true orderings (e.g., proportion of correctly ordered pairs). 
This tests whether the model can infer temporal progression from unordered event content alone.

\paragraph{Survival Analysis for Mortality Time.}  
We include a classical survival analysis task to model time until death. 
Many case reports specify whether the patient died and, if so, when (e.g., "the patient died on hospital day 10"), which enables us to define a time-to-event outcome. 
For this task, we evaluate models at predefined cut-off times -- specifically at 0 hours (admission), 24 hours (1 day), and 168 hours (1 week) -- and use the extracted event history up to each cut-off as input. 
A survival model is trained to predict the probability of survival over time beyond each cut-off. 
We evaluate using the time-dependent concordance index to measure alignment of the predicted survival times with actual outcomes.

\vspace{-.5em}
\subsection{Modeling Approaches}
\vspace{-.25em}
To evaluate performance on the event forecasting and survival prediction tasks, we implement five modeling paradigms: (i) decoder-only large language models (LLMs) with fine-tuned heads, (ii) prompted LLMs without gradient updates (zero- or few-shot), (iii) encoder-only models with task-specific fine-tuned heads, (iv) encoder-masking models with fine-tuning, and (v) encoder-masking models in zero-shot settings. In Appendix~\ref{apd:encoder-forecasting} and ~\ref{apd:encoder-survival}, we provide detailed explanations of implementation details.

\paragraph{(i) Fine-tuned LLMs.}  
We apply instruction-tuned decoder-only models from the \textit{LLaMA} and \textit{DeepSeek} families: \textit{Llama-3.3-70B-Instruct}, \textit{Llama-3.1-8B-Instruct}, and their distilled variants \citep{grattafiori2024llama, liu2024deepseek}. We also include open-source models with documented training corpora (\textit{OLMO-32B-Instruct} \citep{olmo20242}, \textit{RedPajama-INCITE-7B-Instruct} \citep{weber2024redpajama}) to address potential data leakage concerns--notably, RedPajama explicitly excludes PubMed from its training data. Additionally, we evaluate a medically fine-tuned decoder model (\textit{MediPhi-PubMed} \citep{corbeil2025modular}) to assess domain-specific benefits.

Each model is paired with a lightweight multilayer perceptron (MLP) head trained for classification or ranking, depending on the task. 
The input consists of a text-formatted prefix of events (e.g., a clinical timeline up to time $t$), optionally accompanied by an instruction template. 
The output layer produces task-specific predictions: binary labels for event occurrence or a permutation over $k$ events for ordering. 
Training is conducted using cross-entropy loss for classification and pairwise ranking loss for ordering.

\paragraph{(ii) Prompted LLMs.}  
In the zero- or few-shot setting, we use the same LLM architectures as above, but without any fine-tuning. 
Instead, we supply structured prompts at inference time to guide the model toward task-specific outputs. 
Each prompt includes: 
(1) a system instruction establishing the model’s role (e.g., "You are an expert physician."), 
(2) a user instruction describing the prediction task, and 
(3) one or more input-output examples in a few-shot format.
Prompt designs are customized for each task and are detailed in Appendix~\ref{apd:prompts}. Model outputs are parsed to extract binary labels or ordered lists, and errors in generation or ambiguity are excluded from evaluation.

\paragraph{(iii) Fine-tuned Encoder-Only Models.}  
We also evaluate a range of encoder-based models trained end-to-end on each task. 
Architectures include general-purpose models (\textit{BERT-base-uncased} \citep{devlin2019bert}, \textit{RoBERTa-base} \citep{liu2019roberta}, \textit{DeBERTa-v3-small} \citep{he2021debertav3}, \textit{ModernBERT-base/large} \citep{warner2024smarter}) and biomedically-pretrained variants (\textit{BioClinical-ModernBERT-base/large} \citep{sounack2025bioclinical}).
For each model, we append a task-specific MLP head for event occurrence or ordering prediction. 
The model input is the same flattened prefix of the event sequence used in other settings, tokenized and formatted according to the respective architecture. 
These models are trained using standard supervised learning objectives and evaluated on the same metrics as other methods: F1 score for event occurrence, and pairwise concordance for temporal ordering. 
For survival analysis, the final hidden states are passed into a time-to-event prediction head and evaluated using the time-dependent concordance index \citep{antolini2005time}.

\paragraph{(iv, v) Encoder-Masking Models.}
We adapt masked language modeling (MLM) for our temporal reasoning tasks. For event occurrence, models predict masked tokens (``\texttt{yes}''/``\texttt{no}'') in prompts such as ``Will [event] happen within 24 hours? \texttt{[MASK]}''. For temporal ordering, models predict ``\texttt{before}''/``\texttt{after}'' tokens in prompts comparing event pairs. We evaluate both fine-tuned variants (trained on task-specific objectives) and zero-shot variants (using pretrained MLM capabilities without additional training). This approach leverages the bidirectional context of encoder models while maintaining interpretable predictions through constrained vocabulary.

\vspace{-.5em}
\subsection{Sensitivity Analyses}
\vspace{-.25em}
We conduct sensitivity analyses to examine how the temporal structure and completeness of input sequences affect forecasting performance. These analyses probe the robustness of models to changes in event ordering and to varying levels of missing historical information:

\paragraph{Time-ordering Strategies.} We evaluate two strategies for ordering clinical events within textual time series inputs. In the \emph{text-ordered} setting, events are presented in the order they appear in the original case report narratives, preserving the narrative structure produced during extraction. In the \emph{time-ordered} setting, events are sorted chronologically by their documented occurrence time, enforcing strict temporal alignment. This comparison isolates the effect of narrative sequencing versus explicit temporal ordering on model performance. For text-ordered inputs, to prevent causal leakage, events occurring after the forecast cutoff time $t$ are masked for encoder-based models and omitted for decoder-based models.

\paragraph{Timestep Dropout.} To assess robustness to incomplete patient histories, we introduce a \emph{timestep dropout rate} (TDR), defined as the proportion of input timesteps randomly removed prior to inference. This procedure simulates varying degrees of missing clinical documentation that may arise in deployment settings. We vary TDR from 0\% (full history) to 90\% (severely truncated history) in increments, masking events independently and uniformly at random. This setup enables controlled evaluation of model performance under partial information, probing the extent to which predictive accuracy and event ordering depend on the completeness of historical context.

\paragraph{} Together, these analyses characterize model robustness to disruptions in temporal coherence and input completeness -- two key challenges in real-world deployment scenarios.